\section{序論}
11月11日朝の観測窓では、生成AIの製品機能と音声インターフェースの拡張が目立った。AnthropicはClaudeで共同作業・表計算・データ可視化の機能を展開し、Fish Audioは聞きながら話せる音声モデルを紹介した。いずれもモデル単体の性能から、実際の作業や会話に組み込む体験へ競争軸が広がる動きとして注目された。

一方、NVIDIA経営陣を巡る未確認の疑惑、OpenAIの売上報道、Anthropicの安全性基準に関する解釈など、事業とガバナンスに関する議論も続いた。研究者による近未来のAIリスク評価やモデル運用の停止機構も話題となった。本号はGrokの収集結果を順序どおりに整理したもので、独立した一次資料の再検証は行っていない。
